\paragraph{Moore--Penrose pseudoinverse as gauge fixing.}

In the present implementation the tensor is reconstructed using the
Moore--Penrose pseudoinverse,

\begin{equation}
\boldsymbol\tau^{(4)}=
\mathbf R_4^{+}(A,B,C)\,\mathbf D^{(4)},
\label{eq:quartic_pseudoinverse}
\end{equation}

where $\mathbf R_4^{+}$ denotes the pseudoinverse of the quartic projection
matrix in the reduced pair--pair tensor space.

This operation does not impose a physical constraint on the quartic tensor.
Instead, it selects a distinguished representative within the affine family
defined in Eq.~\eqref{eq:quartic_affine_family}. The associated gauge-fixing
condition can be written as

\begin{equation}
\tau_{aabb}+\frac{\Sigma-1}{2}\tau_{aacc}
-\frac{\Sigma+1}{2}\tau_{bbcc}=0,
\label{eq:pseudoinverse_gauge_tau}
\end{equation}

or, equivalently,

\begin{equation}
2\tau'_{ab}+(\Sigma-1)\tau'_{ac}-(\Sigma+1)\tau'_{bc}=0.
\label{eq:pseudoinverse_gauge_tauprime}
\end{equation}

Equation~\eqref{eq:pseudoinverse_gauge_tauprime} should therefore be
interpreted as a pseudoinverse gauge condition rather than as an
intrinsic physical constraint on the quartic tensor.

\paragraph{A practical $3+2$ reduction on the pseudoinverse slice.}

The full spectral representation of the reduced quartic tensor is

\begin{equation}
\tau' = R(\alpha,\beta,\gamma)\,
\mathrm{diag}(\lambda_1,\lambda_2,\lambda_3)\,
R(\alpha,\beta,\gamma)^T,
\label{eq:tauprime_spectral_3plus3}
\end{equation}

which provides a natural $3+3$ parametrization in terms of three
eigenvalues and three orientation angles. This representation describes
the complete six-dimensional space of symmetric $3\times3$ tensors.

When the tensor is reconstructed from Watson constants using the
Moore--Penrose pseudoinverse, however, one does not obtain an arbitrary
point of this six-dimensional space, but the specific representative
selected by the pseudoinverse gauge fixing. In the notation introduced
above, this representative satisfies

\begin{equation}
F(\tau')=
2\tau'_{ab}+(\Sigma-1)\tau'_{ac}-(\Sigma+1)\tau'_{bc}=0,
\label{eq:tauprime_pinv_slice}
\end{equation}

with

\begin{equation}
\Sigma=\frac{2A-B-C}{B-C}.
\end{equation}

Equation~\eqref{eq:tauprime_pinv_slice} removes one coordinate from the
full $3+3$ spectral parametrization and therefore defines a practical
$3+2$ description of the pseudoinverse representative.
A convenient choice of independent coordinates is

\[
(\lambda_1,\lambda_2,\lambda_3,\alpha,\beta),
\]

while the remaining Euler angle $\gamma$ is determined implicitly from
the scalar gauge condition

\begin{equation}
F\!\left(
R(\alpha,\beta,\gamma)\,
\mathrm{diag}(\lambda_1,\lambda_2,\lambda_3)\,
R(\alpha,\beta,\gamma)^T
\right)=0.
\label{eq:gamma_from_gauge}
\end{equation}

This $3+2$ representation should not be interpreted as a global,
representation-independent parametrization of all quartic tensors.
Rather, it defines a coordinate chart on the five-dimensional
pseudoinverse slice selected in the inverse reconstruction problem.
The $3+3$ spectral form therefore remains the fundamental tensorial
description, whereas the $3+2$ reduction provides a convenient
parametrization of the specific tensor representative used in the
present implementation.

\paragraph{Representation transforms.}

Within this tensorial framework, a change of principal-axis
representation is performed through three steps:

\begin{enumerate}
\item reconstruct the pair--pair tensor $\boldsymbol\tau^{(4)}$ from
Watson's constants by pseudoinverse reconstruction;

\item permute the tensor according to the corresponding relabeling of
the principal axes;

\item reproject the permuted tensor onto Watson's quartic constants by
the unique forward map of
Eqs.~\eqref{eq:T_from_tauprime}--\eqref{eq:watson_S_quartic}.
\end{enumerate}

This procedure replaces the explicit analytical transformation matrices
usually written for specific operator parametrizations by a tensor-based
algorithm defined directly in the pair--pair tensor space.

\paragraph{Role of the spectral description.}

Equation~\eqref{eq:tauprime_spectral} shows that the quartic tensor
admits a natural $3+3$ parametrization in terms of three eigenvalues and
three orientation coordinates. Since Watson's quartic Hamiltonian
contains only five independent parameters, the six-dimensional tensor
space must necessarily be reduced when passing to the spectroscopic
constants.

Importantly, this reduction does not correspond to discarding one Euler
angle. Instead, the five-dimensional spectroscopic space is obtained as
the quotient of the six-dimensional tensor space by the
one-dimensional affine gauge family generated by the null tensor $N'$.

From this viewpoint the spectroscopic constants describe equivalence
classes of quartic tensors, whereas the tensor $\tau'$ itself retains
the full spatial information prior to the gauge choice associated with
inverse reconstruction. This distinction is useful both conceptually
and algorithmically, since it separates the intrinsic tensorial
structure of centrifugal distortion from the specific parametrization
chosen in Watson's Hamiltonian.
